% !TeX root = ../../Book4.tex
% 纲要标题：## 第18章　导出范畴的构造
% 纲要标签：b4:ch:17
% 纲要位置：计划纲要.md，第 3682 行。

\chapter{导出范畴的构造}
\label{b4:ch:17}
\BookChapterNumber{b4:ch:17}

\begin{chapterabstract}
本章把拟同构在同伦范畴中形式地变为可逆，建立屋顶演算与导出范畴的三角结构，再用有界投射、内射模型计算导出态射与 Ext。
\end{chapterabstract}

\begin{learninggoals}
\begin{itemize}
\item 用泛性质区别导出范畴与仅列出上同调对象的做法。
\item 通分、合成和比较屋顶，说明拟同构的消去条件。
\item 证明导出范畴的三角结构及短正合列所产生的好三角。
\item 区分项有界、上同调有界和特殊消解的有界方向。
\item 用投射或内射模型计算导出态射，并核对 Ext 的次数与符号。
\end{itemize}
\end{learninggoals}

自由模的两项复形 \(P=(\mathbb Z\xrightarrow{2}\mathbb Z)\) 向只在零度放置 \(\mathbb Z/2\mathbb Z\) 的复形有一个自然的链映射：末项取商，另一项送到零。它在同调上是同构，却没有逐项的逆映射。若计算结果本来就不应取决于选用 \(P\) 还是直接使用商模，便要允许这样的链映射具有逆，而不能要求逆仍由逐项映射给出。导出范畴把这类拟同构正式变为可逆；「屋顶」表示记录了先更换模型、再施行态射的过程。接下来的任务是证明这种新态射仍与锥和三角结构相容。

\ProseParagraphBreak
本章使用\chapref{b4:ch:16}的三角公理和 Hom 正合性。\secref{b4:sec:1704}的模型另外调用\chapref{b4:ch:07}的投射、内射消解，\chapref{b4:ch:08}的 Ext，\cref{b4:thm:ce-resolution}中有下界复形的 Cartan--Eilenberg 消解，以及\cref{b4:cor:double-complex-comparison}的双复形比较。前三节建立导出范畴及其三角结构；\secref{b4:sec:1704}进一步给出计算态射所需的有界模型。

\ProseParagraphBreak
局部化与屋顶可对照 Weibel\cite[Section 10.3; Proposition 10.4.1]{Weibel1994HomologicalAlgebra}；有界消解模型见该书\cite[Corollary 10.4.7; Theorem 10.4.8]{Weibel1994HomologicalAlgebra}。本章同时说明一般交换范畴的集合大小约定，并在模范畴中给出局部小性的具体证明。

% !TeX root = ../../Book4.tex
% 纲要标题：### 18.1 局部化
% 纲要标签：b4:sec:1701
% 纲要位置：计划纲要.md，第 3684 行。

\section{局部化}
\label{b4:sec:1701}

拟同构保存同调，却未必已有同伦逆。将全部拟同构正式变为可逆态射，能允许更换消解，同时仍保留态射的复合及其泛性质。


% 纲要内容（保留纲要核验项目）：
%
% * 把拟同构变成可逆态射
% * 局部化的泛性质
% * 集合大小与构造假设
%

\subsection{拟同构的可逆化}
\label{b4:subsec:1701:01}

链同伦等价必为拟同构，但反向不成立。消解 \(P\to M\) 的作用恰在于：虽然它可能不是同伦等价，上同调信息却与 \(M\) 完全相同。若希望既允许换消解又保持态射的复合，就应把所有拟同构都当作可逆态射。

\begin{definition}\label{b4:def:derived-category}
设 \(\mathcal A\) 为局部小交换范畴，\(S\) 为 \(K(\mathcal A)\) 中的拟同构，即在各次上同调上诱导同构的态射类。在适当的集合大小约定下，取局部化
\[
 D(\mathcal A)=K(\mathcal A)[S^{-1}].
\]
所得范畴称为 \(\mathcal A\) 的\term[daochu-fanchou]{导出范畴}[derived category]，其局部化函子记为 \(Q:K(\mathcal A)\to D(\mathcal A)\)。这里的逆是新范畴中的逆，不要求拟同构原本具有复形映射意义的逆。
\end{definition}
\symidx{\(D(\mathcal A)\)：复形的导出范畴}

一个零调复形在 \(D(\mathcal A)\) 中成为零，因为它到零复形的映射是拟同构。反过来，各次上同调函子能经 \(Q\) 分解，所以在 \(D(\mathcal A)\) 中为零的对象必零调。因此，复形 \(C\) 在 \(K(\mathcal A)\) 中为零对象当且仅当它可缩，在 \(D(\mathcal A)\) 中为零对象当且仅当它零调；零调但不可缩的复形也在局部化后与零复形同构。

\ProseParagraphBreak
例如，将正合列
\[
 0\longrightarrow\mathbb Z\xrightarrow{2}\mathbb Z
 \xrightarrow{q}\mathbb Z/2\mathbb Z\longrightarrow0
\]
的三个非零项依次放在次数零、一、二，得到零调复形 \(C\)。若它可缩，同伦等式在次数二要求 \(q h^2=1_{\mathbb Z/2\mathbb Z}\)，其中 \(h^2:\mathbb Z/2\mathbb Z\to\mathbb Z\)。但这样的同态只能为零，不可能给出 \(q\) 的右逆。因此 \(C\) 在同伦范畴中非零，在导出范畴中却为零。

\subsection{局部化的泛性质}
\label{b4:subsec:1701:02}

\begin{definition}\label{b4:def:category-localization}
设 \(\mathcal C\) 为范畴，\(S\) 为一类态射，给定函子 \(Q:\mathcal C\to\mathcal C[S^{-1}]\)。若 \(Q(s)\) 对每个 \(s\in S\) 都可逆，并且对每个范畴 \(\mathcal E\)，任何将 \(S\) 送为同构的函子 \(F:\mathcal C\to\mathcal E\) 都唯一写成 \(F=\overline FQ\)，自然变换也唯一经此分解，则称 \(Q\) 为 \(\mathcal C\) 在 \(S\) 处的\term[fanchou-jubuhua]{范畴局部化}[localization of a category]。采用对象不变的标准构造时，这里的唯一性是严格的；更换等价模型后则为唯一到自然同构。
\end{definition}
\symidx{\(\mathcal C[S^{-1}]\)：在指定态射类处的范畴局部化}

小范畴的一个直接构造，是在原有箭头之外为每个 \(s\in S\) 添一个反向符号 \(s^{-1}\)，取有限可复合箭头串，再按原范畴的复合关系以及 \(ss^{-1}=1,s^{-1}s=1\) 取商。复合由连接箭头串给出。函子 \(F\) 在新符号上只能取 \(F(s)^{-1}\)，因而唯一延拓；自然变换在新符号上的自然性由在 \(s\) 上的自然性乘逆得到。这也证明了自然变换的分解性质。

\ProseParagraphBreak
形式逆不必来自原范畴中的映射。取位于次数负一、零的 \(P=(\mathbb Z\xrightarrow{2}\mathbb Z)\)，以及拟同构 \(s:P\to(\mathbb Z/2\mathbb Z)[0]\)。任何反向复形映射 \(a:(\mathbb Z/2\mathbb Z)[0]\to P\) 在次数零都是 \(\mathbb Z/2\mathbb Z\to\mathbb Z\) 的零同态，故不能满足 \(sa=1\)，也不能成为 \(s\) 的同伦逆。局部化仍使 \(Q(s)^{-1}\) 存在；下一节的屋顶将通过把源对象换成 \(P\)，而非寻找这样的 \(a\)，来表示这个逆。

\ProseParagraphBreak
这种构造证明存在，却不便计算。下一节证明：同伦范畴中的拟同构允许把任意箭头串缩成一个分母在左的屋顶。这里不能先假定任意范畴都有这样的分式形式；它需要专门的交换与消去条件。

\ProseParagraphBreak
还可直接从复形范畴把拟同构倒过来。两种做法等价：拟同构包含链同伦等价，故\cref{b4:prop:homotopy-universal}迫使任何这样的函子先经 \(K(\mathcal A)\) 分解；之后再用局部化的泛性质即可。

\subsection{集合大小与构造假设}
\label{b4:subsec:1701:03}

若 \(\mathcal A\) 的对象不是一个集合，前述箭头串也可能形成真类，不能仅凭原范畴局部小就断言局部化仍局部小。本书对一般交换范畴使用固定的较大集合范围来作局部化；讨论 Hom 群时，要么假定局部化在所用范围内局部小，要么使用下面模范畴的大小控制，或\secref{b4:sec:1704}的有界消解模型。任意交换范畴的无界情形还需要额外的局部小性证明。

\begin{proposition}\label{b4:prop:module-derived-size}
设 \(R\) 为含幺环，\(X\) 为左 \(R\)-模复形。存在一集合的拟同构 \(X_i\to X\)，使每个拟同构 \(s:Y\to X\) 都能由某个 \(X_i\to X\) 经拟同构 \(X_i\to Y\) 细化。因此在允许屋顶演算后，\(D(R\text{-}\mathbf{Mod})\) 的每个态射类确实组成集合。
\end{proposition}
\begin{proof}
取无限基数 \(\kappa\) 至少为 \(|R|\)、\(\aleph_0\) 及所有 \(|X^n|\) 的上界。给定 \(s\)，为 \(X\) 每次上同调的每个类选一个 \(Y\) 中余圈，其像代表该类。所选元素生成一个各次大小至多 \(\kappa\) 的子复形 \(Y_0\subseteq Y\)。这一步保证 \(H(Y_0)\to H(X)\) 满，但缩小到子复形后，一个在 \(Y\) 中为余边的余圈未必仍在 \(Y_0\) 中为余边，所以还需消去这个映射的核。

归纳构造 \(Y_j\)：对 \(Y_j\) 中每个余圈 \(z\)，若 \(s(z)\) 在 \(X\) 中为余边，则 \(s\) 的上同调单射性保证 \(z=\mathrm{d}y\) 对某个 \(y\in Y\) 成立。把这些 \(y\) 加入并生成子复形，得 \(Y_{j+1}\)。新加入的元素还可能产生新的余圈，因此要重复这个过程。每一步至多加入 \(\kappa\) 个元素；取可数并 \(Y'=\bigcup_jY_j\)，其大小仍至多 \(\kappa\)。起初的选择保证 \(H(Y')\to H(X)\) 满。任一落入核的余圈 \(z\in Y'\) 已属于某个 \(Y_j\)，它在 \(Y_{j+1}\) 中成为余边，所以该映射也单。于是 \(Y'\to X\) 是拟同构；它是 \(Y'\to Y\to X\) 的复合，后二者中的 \(Y\to X\) 也为拟同构，故 \(Y'\to Y\) 同样为拟同构。

所有大小至多 \(\kappa\) 的模复形可以在固定集合上编码：每次的底集合取为一个固定的 \(\kappa\) 元集合的子集，模运算和微分都是这些集合上的函数，可能的选择组成集合。它们到 \(X\) 的复形映射也组成集合。取其代表即得所需的一集合分母。每个屋顶 \(X\leftarrow Y\to Z\) 都可换成以某个 \(X_i\) 为中间对象的屋顶，而 \(\bigcup_i\operatorname{Hom}_{K}(X_i,Z)\) 为集合，其商仍为集合。
\end{proof}

这里用 \(\le\kappa\) 控制大小，避免了“可数次合并以后仍严格小于某个基数”的不当推断。相关大小问题可对照 Weibel\cite[Set-Theoretic Considerations 10.3.6; Proposition 10.4.4]{Weibel1994HomologicalAlgebra}。

% !TeX root = ../../Book4.tex
% 纲要标题：### 18.2 屋顶与分式
% 纲要标签：b4:sec:1702
% 纲要位置：计划纲要.md，第 3690 行。

\section{屋顶与分式}
\label{b4:sec:1702}

导出态射可以先沿拟同构更换源对象，再作复形映射。屋顶的共同细化既给出相等判据，也使复合、加法与三角结构可以严格构造。


% 纲要内容（保留纲要核验项目）：
%
% * 屋顶图表示态射
% * 适当的分式演算
% * 同伦范畴到导出范畴
%

\subsection{态射的屋顶图表示}
\label{b4:subsec:1702:01}

\begin{definition}\label{b4:def:derived-roof}
设 \(\mathcal A\) 为局部小交换范畴，在固定的集合大小约定下记局部化函子为 \(Q:K(\mathcal A)\to D(\mathcal A)\)。\(K(\mathcal A)\) 中的图式
\[
 X\xleftarrow{s}U\xrightarrow{f}Y,\qquad s\;\text{为拟同构},
\]
称为一个\term[wuding]{屋顶}[roof]，表示 \(Q(f)Q(s)^{-1}:X\to Y\)。给定另一个屋顶 \(X\xleftarrow{s'}U'\xrightarrow{f'}Y\)，若存在对象 \(V\) 及拟同构 \(a:V\to U\)、\(b:V\to U'\)，使 \(sa=s'b\)、\(fa=f'b\)，则称两个屋顶有共同细化；所有等式都在同伦范畴中理解。
\end{definition}

共同细化把两个屋顶放在同一张图中：
\[
\begin{tikzcd}[column sep=3.2em,row sep=2.2em]
&U\arrow[dl,"s\;(\sim)"']\arrow[dr,"f"]&\\
X&V\arrow[u,"a\;(\sim)"]\arrow[d,"b\;(\sim)"']&Y\\
&U'\arrow[ul,"s'\;(\sim)"]\arrow[ur,"f'"']&
\end{tikzcd}
\]
左侧要求 \(sa=s'b\)，右侧要求 \(fa=f'b\)。这里 \(a,b\) 都是拟同构；也可以先要求共同左分母 \(sa=s'b\) 为拟同构，再由二取三性质推出这一点。符号 \(Q(s)^{-1}\) 是局部化以后得到的逆，图式本身只使用同伦范畴中实际存在的箭头，没有要求在那里选择 \(s\) 的逆映射。

\ProseParagraphBreak
在模消解的典型例子中，\(P\to M\) 是左分母，而 \(P\to N[n]\) 记录一个 Ext 余圈。屋顶因此不是额外的画图规则，而是“先换成可计算的模型，再作映射”的正式表达。接下来证明共同细化足以判别相等。

\subsection{分式演算}
\label{b4:subsec:1702:02}

屋顶的复合需要把指向同一对象的两张箭头通分；比较不同的通分选择，则需要消去两张箭头的差。下面第一项构造等式 \(ft=sg\)，使复合所需的方块交换；第二项从 \(sf=sg\) 出发，允许进一步前复合一个拟同构后得到 \(ft=gt\)，用来保证计算不依赖选择。第三、第四项分别是后复合方向的通分和消去，它们给出分母在右的分式演算。

\begin{lemma}\label{b4:lem:quasi-ore}
设 \(\mathcal A\) 为局部小交换范畴。\(K(\mathcal A)\) 中的拟同构包含恒等态射，对复合和平移封闭，并满足二取三性质。它们还满足以下四个条件。
\begin{enumerate}
\item 给定 \(f:X\to Y\) 和拟同构 \(s:Z\to Y\)，存在拟同构 \(t:W\to X\) 和 \(g:W\to Z\)，使 \(ft=sg\)。
\item 若拟同构 \(s:Y\to Y'\) 满足 \(sf=sg\)，其中 \(f,g:X\to Y\)，则存在拟同构 \(t:X'\to X\)，使 \(ft=gt\)。
\item 给定 \(f:X\to Y\) 和拟同构 \(s:X\to Z\)，存在拟同构 \(t:Y\to W\) 和 \(g:Z\to W\)，使 \(tf=gs\)。
\item 若拟同构 \(s:X'\to X\) 满足 \(fs=gs\)，其中 \(f,g:X\to Y\)，则存在拟同构 \(t:Y\to Y'\)，使 \(tf=tg\)。
\end{enumerate}
\end{lemma}
\begin{proof}
恒等态射、复合、平移及二取三性质的结论直接对各次上同调核对。对第一项，补全 \(Z\xrightarrow{s}Y\to C\to Z[1]\)。因 \(s\) 为拟同构，\(C\) 零调。再把复合 \(X\xrightarrow{f}Y\to C\) 补成旋转后的好三角 \(W\xrightarrow{t}X\to C\to W[1]\)。长正合列说明 \(t\) 是拟同构；对这两个三角使用 (TR3)，即可得到 \(g\) 及 \(ft=sg\)。

对第二项令 \(h=f-g\)。补全 \(s\) 的锥三角，记 \(C=\operatorname{Cone}(s)\)。由 \(sh=0\) 和\cref{b4:lem:pretriangle-hom}的协变 Hom 正合性，\(h\) 经 \(C[-1]\) 分解为 \(X\xrightarrow{a}C[-1]\to Y\)。这里 \(C[-1]\) 虽然零调，却未必是 \(K(\mathcal A)\) 中的零对象，不能仅由这一分解断言 \(h=0\)。把 \(a\) 补成好三角
\[
 X'\xrightarrow{t}X\xrightarrow{a}C[-1]\longrightarrow X'[1].
\]
由\cref{b4:prop:cohomology-on-homotopy}产生的上同调长正合列及 \(C[-1]\) 零调，\(t\) 是拟同构。相邻复合为零给出 \(at=0\)，所以 \(ht=0\)，即 \(ft=gt\)。零调锥在这里提供的是进一步拟同构后的消去，而非原同伦范畴中的直接消去。

对第三项，补全 \(X\xrightarrow{s}Z\to C\to X[1]\)，其中 \(C\) 零调。逆向旋转后记首箭头为 \(u:C[-1]\to X\)，再把 \(fu:C[-1]\to Y\) 补成好三角
\[
 C[-1]\xrightarrow{fu}Y\xrightarrow{t}W\to C.
\]
对逆向旋转的三角和这个三角应用 (TR3)，前两个分量取 \(1_{C[-1]},f\)，得到 \(g:Z\to W\)，并有 \(tf=gs\)。长正合列及 \(C\) 零调说明 \(t\) 是拟同构。

对第四项令 \(h=f-g\)，补全 \(X'\xrightarrow{s}X\xrightarrow{q}C\to X'[1]\)。由 \(hs=0\) 和反变 Hom 正合性，存在 \(b:C\to Y\) 使 \(h=bq\)。把 \(b\) 补成好三角 \(C\xrightarrow{b}Y\xrightarrow{t}Y'\to C[1]\)。\(C\) 零调，故 \(t\) 是拟同构，而 \(th=tbq=0\)。
\end{proof}

\begin{theorem}\label{b4:thm:roof-calculus}
设 \(\mathcal A\) 为局部小交换范畴，并采用\subsecref{b4:subsec:1701:03}的集合大小约定，记局部化函子为 \(Q:K(\mathcal A)\to D(\mathcal A)\)。\(D(\mathcal A)\) 中每个态射可由屋顶表示，两个屋顶表示同一态射当且仅当它们有共同细化。特别地，对 \(K(\mathcal A)\) 中的 \(f,g:X\to Y\)，\(Q(f)=Q(g)\) 当且仅当存在拟同构 \(t:X'\to X\) 使 \(ft=gt\)；也可用后复合的拟同构 \(Y\to Y'\) 检验。
\end{theorem}
\begin{proof}
先用共同细化定义屋顶的等价。自反、对称显然。对传递性，把两次细化到共同中间屋顶的分母用\cref{b4:lem:quasi-ore}第一项再作共同细化，二取三保证新添的箭头仍为拟同构；复合之后，左右两条边同时相等。这证明它确为等价关系。

复合 \((s,f):X\to Y\) 与 \((t,g):Y\to Z\) 时，对 \(f:U\to Y\) 和 \(t:V\to Y\) 取 \(a:W\to U\)、\(b:W\to V\)，其中 \(a\) 拟同构且 \(fa=tb\)。拼接图为
\[
\begin{tikzcd}[column sep=1.8em,row sep=2.1em]
&&W\arrow[dl,"a\;(\sim)"']\arrow[dr,"b"]&&\\
&U\arrow[dl,"s\;(\sim)"']\arrow[dr,"f"]&&V\arrow[dl,"t\;(\sim)"']\arrow[dr,"g"]&\\
X&&Y&&Z.
\end{tikzcd}
\]
由 \(fa=tb\)，在局部化中有
\[
 Q(t)^{-1}Q(f)=Q(b)Q(a)^{-1}.
\]
因此两个屋顶的复合为
\[
 Q(g)Q(t)^{-1}Q(f)Q(s)^{-1}
 =Q(gb)Q(sa)^{-1},
\]
正是沿图的外侧两条路径得到的屋顶 \((sa,gb)\)。图中标了 \(\sim\) 的 \(s,t,a\) 必须为拟同构，才能使用上述逆；辅助箭头 \(b\) 无须为拟同构。若另一选择为 \((a',b')\)，先共同细化拟同构 \(a,a'\)，得到 \(ac=a'c'\)。于是 \(tbc=tb'c'\)，用引理第二项再前复合一个拟同构，便使 \(bc=b'c'\) 也成立。因此两复合屋顶等价。更换原屋顶的代表时，把其共同细化与这一计算合并，同样得到等价的复合。

三重屋顶的两种括号分别使用有限个这样的交换方块。先将它们涉及的分母逐一共同细化，再用消去条件使所有指向同一对象的两条有限路径相等，就得到一个同时计算两种括号的屋顶；其左边是各分母的复合，右边是各分子的复合。因此结合律归结为 \(K(\mathcal A)\) 的结合律。这里每次只处理三重复合中的有限个方块，不涉及无限图式的同时提升。恒等屋顶为 \((1,1)\)，单位律由取恒等交换方块得到。

函子 \(Q\) 把 \(f\) 送到 \((1,f)\)。若 \(s\) 拟同构，\((s,1)\) 是 \(Q(s)\) 的逆。任意倒转拟同构的函子 \(F\) 必须把屋顶送到 \(F(f)F(s)^{-1}\)；共同细化保证此式与代表无关，复合公式保证它保持复合。这就给出唯一的分解函子 \(\overline F\)。

还需核对自然变换的分解。设 \(F,G:K(\mathcal A)\to\mathcal E\) 都把拟同构送为同构，\(\eta:F\Rightarrow G\) 为自然变换。屋顶范畴的对象不变，所以只能取 \(\overline\eta_X=\eta_X\)。对屋顶 \(X\xleftarrow{s}U\xrightarrow{f}Y\)，由 \(\eta\) 的自然性有
\[
 \eta_YF(f)=G(f)\eta_U,\qquad
 \eta_XF(s)=G(s)\eta_U.
\]
因 \(F(s),G(s)\) 可逆，得到
\[
 \eta_Y\overline F(s,f)
 =G(f)\eta_UF(s)^{-1}
 =G(f)G(s)^{-1}\eta_X
 =\overline G(s,f)\eta_X.
\]
因此这组分量确实给出唯一的自然变换 \(\overline\eta:\overline F\Rightarrow\overline G\)。屋顶范畴满足局部化的全部泛性质，正是 \(D(\mathcal A)\)。最后把两个普通态射看作分母恒等的屋顶，共同细化条件便是所写的相等判据；后复合的判据由引理第三、第四项得到。
\end{proof}

对两个屋顶 \(X\xleftarrow{s}U\xrightarrow{f}Y\)、\(X\xleftarrow{t}V\xrightarrow{g}Y\)，用通分条件取 \(a:W\to U\)、\(b:W\to V\)，使 \(sa=tb\) 为拟同构。二取三性质保证 \(a,b\) 也都是拟同构。两张右箭头此时有同一源对象，才可相加，所得屋顶为
\[
 X\xleftarrow{sa}W\xrightarrow{fa+gb}Y.
\]
同分母时，这就是 \((s,f)+(s,g)=(s,f+g)\)。不同的通分选择可进一步共同细化，并用消去条件使相应右箭头相等；预复合又保持加法，故结果与选择无关。复合公式及原来的双线性给出新复合的双线性。零对象和有限双积由 \(K(\mathcal A)\) 保留下来。因此 \(D(\mathcal A)\) 是加性范畴，\(Q\) 为加性函子。屋顶演算的一般形式可参阅 Weibel\cite[Theorem 10.3.7; Corollaries 10.3.8--10.3.11]{Weibel1994HomologicalAlgebra}。

\subsection{从同伦范畴到导出范畴}
\label{b4:subsec:1702:03}

\begin{theorem}\label{b4:thm:derived-triangulated}
设 \(\mathcal A\) 为局部小交换范畴，并采用\subsecref{b4:subsec:1701:03}的集合大小约定，记局部化函子为 \(Q:K(\mathcal A)\to D(\mathcal A)\)。在 \(D(\mathcal A)\) 中，以平移诱导的自等价为平移，以同构于 \(K(\mathcal A)\) 中好三角之像的三角为好三角，则 \(D(\mathcal A)\) 是三角范畴，\(Q\) 保持平移和好三角。
\end{theorem}
\begin{proof}
平移保持拟同构，所以经局部化诱导自等价。(TR1) 中的恒等三角和对同构封闭性已成立。任意屋顶 \(X\xleftarrow{s}U\xrightarrow{f}Y\) 的左箭头在 \(D\) 中可逆，因此 \(f\) 的锥三角运送后给出该屋顶的好三角。(TR2) 由 \(K\) 中的旋转直接得到。

核对 (TR3) 时，先把两个三角换成 \(K\) 中标准三角的像，首箭头记为 \(f:X\to Y\)、\(f':X'\to Y'\)。给定导出范畴中的交换方块 \((\alpha,\beta)\)，分别取屋顶
\(X\xleftarrow{s}S\xrightarrow{a}X'\)、\(Y\xleftarrow{t}T\xrightarrow{b}Y'\)。沿拟同构 \(t\) 通分 \(fs\)，得到拟同构 \(u:U\to S\) 和 \(k:U\to T\)，满足 \(t k=fsu\)。方块在 \(D\) 中交换，故 \(Q(f'au)=Q(bk)\)。由屋顶相等判据，再取拟同构 \(v:V\to U\)，使 \(f'auv=bkv\) 在 \(K\) 中成立。

上述构造给出 \(K\) 中的两个交换方块，它们共享上方箭头 \(kv:V\to T\)，其交换等式为
\[
 t(kv)=f(suv),\qquad b(kv)=f'(auv).
\]
把 \(kv\) 补成一条标准锥三角，再分别用 \(K\) 中的 (TR3) 映向 \(f\)、\(f'\) 的三角：前一对分量为 \((suv,t)\)，后一对为 \((auv,b)\)。这些方块的等式在 \(K\) 中成立，所选的复形映射代表可以只同伦交换；前章的锥映射公式正是用同伦修正它们。前一补全的第三个分量是拟同构，因为前两个分量是拟同构，对各次上同调长正合列使用\cref{b4:thm:five-lemma}即可。因此这一三角态射在 \(D\) 中成为同构；先取其逆，再复合后一三角态射，就得到原方块的补全。这证明 (TR3)，并使\cref{b4:lem:pretriangle-hom}现在也可在 \(D\) 中使用。

对 (TR4)，设 \(\alpha:X\to Y\)、\(\beta:Y\to Z\) 为两个导出态射。要对它们实际构造锥，需先用屋顶通分，将这一对导出态射换成 \(K\) 中可复合的复形映射。先取 \(\beta\) 的屋顶 \(Y\xleftarrow{t}T\to Z\)，再取 \(\alpha\) 的屋顶 \(X\xleftarrow{s}S\to Y\)，并沿 \(t\) 通分右边。得到 \(K\) 中一串 \(U\to T\to Z\)，以及拟同构 \(U\to X\)、\(T\to Y\)，其局部化像正好代表原来的可复合对。在 \(K\) 中对此串应用\cref{b4:thm:homotopy-triangulated}证明中的显式八面体，再施加 \(Q\) 并沿上述两个同构运送，便得到补全 \(\alpha,\beta,\beta\alpha\) 的三个好三角及它们的八面体。

这样得到的是由所选复形映射构造的三角，而 (TR4) 还要求八面体适配任意事先指定的三个好三角。对每个首箭头 \(f\in\{\alpha,\beta,\beta\alpha\}\)，在所得三角与指定三角的前两项取恒等态射，由已证的 (TR3) 补出第三分量 \(c_f:C_f\to C_f'\)。\cref{b4:lem:pretriangle-hom}只需前三条公理，故可在这里应用：前两个分量为同构，所以 \(c_f\) 也可逆。把所得八面体在第三对象之间的箭头记为 \(a:C_\alpha\to C_{\beta\alpha}\)、\(b:C_{\beta\alpha}\to C_\beta\)、\(\gamma:C_\beta\to C_\alpha[1]\)，定义运送后的箭头为
\[
 a'=c_{\beta\alpha}a c_\alpha^{-1},\qquad
 b'=c_\beta b c_{\beta\alpha}^{-1},\qquad
 \gamma'=c_\alpha[1]\gamma c_\beta^{-1}.
\]
各个 \(c_f\) 都与对应三角的后两条箭头相容，故八面体的四个相容等式保持成立；最后一个面经 \((c_\alpha,c_{\beta\alpha},c_\beta)\) 成为同构的好三角。这便满足针对任意指定三角的 (TR4)。
\end{proof}

\begin{proposition}\label{b4:prop:derived-short-exact-triangle}
设 \(\mathcal A\) 为局部小交换范畴，\(0\to A\xrightarrow{i}B\xrightarrow{q}C\to0\) 是 \(\mathcal A\) 中上链复形的短正合列，不要求逐次分裂。记 \(Q:K(\mathcal A)\to D(\mathcal A)\) 为局部化函子，\(p_i:\operatorname{Cone}(i)\to A[1]\) 为锥的投影，令 \(r:\operatorname{Cone}(i)\to C\) 为 \(r(b,a)=q(b)\)。则 \(r\) 是拟同构，且
\[
 A\xrightarrow{i}B\xrightarrow{q}C\xrightarrow{\delta}A[1],
 \qquad\delta=Q(p_i)Q(r)^{-1},
\]
是 \(D(\mathcal A)\) 中的好三角。
\end{proposition}
\begin{proof}
\(r\) 是逐次满的复形映射，因为 \(q\) 逐次满且 \(qi=0\)。利用 \(i^n\) 识别 \(\ker q^n\) 与 \(A^n\)，有
\[
 (\ker r)^n=\ker q^n\oplus A^{n+1}
 \cong A^n\oplus A^{n+1}.
\]
在此识别下，锥的微分为
\[
 \mathrm{d}_{\ker r}^n(a,a')
 =(\mathrm{d}_A^na+a',-\mathrm{d}_A^{n+1}a'),
\]
恰好是 \(\operatorname{Cone}(1_A)\) 的微分。该锥由\cref{b4:prop:cone-sign-compatibility}可缩，因而零调。对
\(0\to\ker r\to\operatorname{Cone}(i)\xrightarrow{r}C\to0\)
应用\cref{b4:thm:cohomology-long-exact}，得到每个 \(H^n(r)\) 为同构，故 \(r\) 拟同构。又因 \(r\) 与锥的包含复合等于 \(q\)，把 \(i\) 的标准锥三角经 \(Q(r)\) 运送，即得到所述好三角及最后箭头 \(Q(p_i)Q(r)^{-1}\)。这只使用 \(r\) 的拟同构性，无须原短正合列具有复形截面。
\end{proof}

若原短正合列逐次分裂，选 \(B^n=A^n\oplus C^n\)，并写
\[
 \mathrm{d}_B=\begin{pmatrix}\mathrm{d}_A&t\\0&\mathrm{d}_C\end{pmatrix}.
\]
由 \(\mathrm{d}_B^2=0\) 得 \(\mathrm{d}_At+t\mathrm{d}_C=0\)。前章\cref{b4:prop:degreewise-split-triangle}中的复形映射截面，在这里记为
\[
 \sigma^n:C^n\longrightarrow\operatorname{Cone}(i)^n,
 \qquad \sigma^n(c)=((0,c),-t^n(c)).
\]
它的微分满足
\[
 \mathrm{d}_{\operatorname{Cone}(i)}\sigma(c)
 =((0,\mathrm{d}_Cc),\mathrm{d}_At(c))
 =\sigma(\mathrm{d}_Cc),
\]
且 \(r\sigma=1_C\)。局部化后 \(Q(r)\) 可逆，因而 \(Q(\sigma)=Q(r)^{-1}\)，于是
\[
 \delta=Q(p_i)Q(r)^{-1}=Q(p_i\sigma)=Q(-t).
\]
负号来自 \(\sigma\) 的第二分量，与前章的正投影约定一致。在模的情形，普通上同调连接同态由余圈 \(c\) 的提升 \((0,c)\) 及其微分 \((t(c),0)\) 计算，给出 \([c]\mapsto[t(c)]\)；这里三角连接箭头在上同调上给出 \([c]\mapsto-[t(c)]\)。一般交换范畴中的比较已由\cref{b4:prop:degreewise-split-triangle}之后的余圈子对象计算给出，同样相差所选锥约定的负号。

\ProseParagraphBreak
此外，\(D\) 中态射为同构，当且仅当其锥为零：用三角的 Hom 正合列，锥为零使该态射对每个 Hom 群诱导同构，从而可逆；反向由三角上同调或同样的 Hom 列得到。对原来的复形映射，这恰好恢复拟同构判据。

% !TeX root = ../../Book4.tex
% 纲要标题：### 18.3 导出范畴中的有界性
% 纲要标签：b4:sec:1703
% 纲要位置：计划纲要.md，第 3696 行。

\section{导出范畴中的有界性}
\label{b4:sec:1703}

复形的项有界与其同调有界并非同一条件。典范截断允许删去正合尾部，却可能失去投射或内射性，因此不能把有界对象与有界特殊消解混同。


% 纲要内容（保留纲要核验项目）：
%
% * D^b、D^+ 与 D^-
% * 项有界和同调有界的区别
% * 模范畴中的基本实例
%

\subsection{\texorpdfstring{\(D^b\)、\(D^+\) 与 \(D^-\)}{Db、D+ 与 D−}}
\label{b4:subsec:1703:01}

\begin{definition}\label{b4:def:bounded-derived}
设 \(\mathcal A\) 为局部小交换范畴，并固定局部化的集合大小约定。分别将 \(K^b(\mathcal A),K^+(\mathcal A),K^-(\mathcal A)\) 中的拟同构倒转，得到\term[youjie-daochu-fanchou]{有界导出范畴}[bounded derived category] \(D^b(\mathcal A)\)、有下界导出范畴 \(D^+(\mathcal A)\) 和有上界导出范畴 \(D^-(\mathcal A)\)。有界在此先指复形的项。
\end{definition}
\symidx{\(D^b(\mathcal A),D^+(\mathcal A),D^-(\mathcal A)\)：有界、有下界与有上界导出范畴}

\begin{proposition}\label{b4:prop:bounded-derived-full}
设 \(\mathcal A\) 为局部小交换范畴，且相关局部化均采用同一集合大小约定。自然函子 \(D^b(\mathcal A),D^+(\mathcal A),D^-(\mathcal A)\to D(\mathcal A)\) 全忠实；它们的本质像分别由上同调两端有界、有下界、有上界的对象组成。三者各为三角子范畴。
\end{proposition}
\begin{proof}
截去尾部时必须保留截口上的余圈与余边条件。例如，\(\mathbb Z\xrightarrow{2}\mathbb Z\xrightarrow{q}\mathbb Z/2\mathbb Z\) 位于次数零、一、二，原复形零调；若直接删去次数大于一的项，余下的 \(\mathbb Z\xrightarrow{2}\mathbb Z\) 却有 \(H^1=\mathbb Z/2\mathbb Z\)。这是因为末项的全部元素都变成余圈。把次数一的项改为 \(\ker q=2\mathbb Z\)，所得 \(\mathbb Z\xrightarrow{2}2\mathbb Z\) 就仍零调。

一般地，复形 \(C\) 的 \(\tau_{\le n}C\) 在 \(n\) 次取 \(\ker \mathrm{d}_C^n\)，较低次数不变，较高次数取零；它自然映入 \(C\)，并保留次数不大于 \(n\) 的上同调。\(\tau_{\ge m}C\) 在 \(m\) 次取 \(\operatorname{coker}\mathrm{d}_C^{m-1}\)，较高次数不变，较低次数取零；\(C\) 自然映到它，并保留次数不小于 \(m\) 的上同调，其余次数的上同调为零。右端取核，是为了保留原来成为余圈的条件；左端取余核，是为了仍将来自被删项的余边模掉。这些断言直接从核、像的定义核对，因而两种截断都保持拟同构。链映射在截口分别限制到核、下降到余核，给出截断后的链映射，并保持恒等与复合。

为了在同伦范畴中截断屋顶及其共同细化，还必须保证同伦的映射经截断后仍同伦。截口取核或余核依赖整个复形，不能仅由逐项加性函子保持同伦的结论代替这一步。设 \(f,g:C\to E\) 及 \(h^i:C^i\to E^{i-1}\) 满足
\[
 f^i-g^i=\mathrm{d}_E^{i-1}h^i+h^{i+1}\mathrm{d}_C^i.
\]
对 \(\tau_{\le n}\)，在 \(i<n\) 保留 \(h^i\)，在 \(i=n\) 将 \(h^n\) 限制到 \(\ker\mathrm{d}_C^n\)，在更高次数取零。截口上 \(h^{n+1}\mathrm{d}_C^n\) 消失，剩下的 \(\mathrm{d}_E^{n-1}h^n\) 取值于 \(\ker\mathrm{d}_E^n\)，所以同伦等式仍成立；在 \(n-1\) 次，限制后的 \(h^n\) 与诱导微分的复合仍为 \(h^n\mathrm{d}_C^{n-1}\)。对 \(\tau_{\ge m}\)，记截口的商映射为
\[
 q_C:C^m\to\operatorname{coker}\mathrm{d}_C^{m-1},
 \qquad q_E:E^m\to\operatorname{coker}\mathrm{d}_E^{m-1}.
\]
在 \(i>m+1\) 保留 \(h^i\)，在 \(i=m+1\) 取 \(q_Eh^{m+1}\)，在 \(i\le m\) 取零。由于 \(q_E\mathrm{d}_E^{m-1}=0\)，截口上的等式成为
\[
 (\tau_{\ge m}f-\tau_{\ge m}g)^m q_C
 =q_Eh^{m+1}\mathrm{d}_C^m,
\]
正是余核上所需的同伦等式；在 \(m+1\) 次，余核诱导的微分与 \(q_Eh^{m+1}\) 复合仍为 \(\mathrm{d}_E^mh^{m+1}\)。因此两种截断都给出 \(K(\mathcal A)\) 上的函子。

若 \(X,Y\) 的项在次数 \(n\) 以上为零，则屋顶 \(X\leftarrow U\to Y\) 可前复合 \(\tau_{\le n}U\to U\)。这里使用左屋顶，是因为截断的自然映射指向 \(U\)。因 \(H^{>n}(U)=0\)，这是拟同构，所以得到有上界的屋顶。若两个有上界的屋顶在 \(D(\mathcal A)\) 中相等，可增大 \(n\)，使这两个屋顶的全部对象在 \(n\) 以上为零，再对它们的共同细化施加 \(\tau_{\le n}\)。截断保持拟同构及 \(K(\mathcal A)\) 中的等式，且不改变两个原屋顶，因此得到有上界的共同细化。这同时证明 \(D^-\) 的满性和忠实性。对 \(D^+\)，截断的自然映射方向为 \(U\to\tau_{\ge m}U\)，因此使用分母在右的屋顶 \(X\to U\leftarrow Y\)，并把两个箭头后复合到截断对象。比较两个屋顶时，取它们全部对象的公共下界 \(m\)，截断共同细化，得到相同的结论。

若 \(X,Y\) 都只在 \([m,n]\) 内有项，先把左屋顶截到 \(\tau_{\le n}U\)，再取它的 \(\tau_{\ge m}\)。指向 \(X,Y\) 的复形映射在次数 \(m\) 都消去来自 \(m-1\) 的像，因此均通过该余核下降，得到有界屋顶；其左边仍为拟同构。比较两个有界屋顶时，选取包含它们全部对象的非零项的公共区间 \([m,n]\)，令 \(T=\tau_{\ge m}\tau_{\le n}\)。若共同细化中的等式为 \(s_1a=s_2b\)、\(f_1a=f_2b\)，施加 \(T\) 后仍有
\[
 T(s_1)T(a)=T(s_2)T(b),
 \qquad T(f_1)T(a)=T(f_2)T(b).
\]
这里 \(T\) 不改变两个原屋顶，而 \(T(a),T(b)\) 仍为拟同构，故所得细化属于 \(K^b(\mathcal A)\)。这证明 \(D^b\) 也全忠实。同样的截断论证在单向有界的屋顶范畴内直接进行，还给出 \(D^b(\mathcal A)\to D^+(\mathcal A),D^-(\mathcal A)\) 的全忠实性。

若 \(H^i(X)=0\) 对所有 \(i\notin[m,n]\) 成立，则有拟同构之字形
\[
 X\xleftarrow{\simeq}\tau_{\le n}X
 \xrightarrow{\simeq}\tau_{\ge m}\tau_{\le n}X.
\]
最右的复形只在 \([m,n]\) 内有项，故 \(X\) 在导出范畴中与项有界的复形同构。只有上界或下界时，分别用 \(\tau_{\le n}X\to X\) 或 \(X\to\tau_{\ge m}X\) 即可。反向由上同调保持拟同构显然。平移和锥保持相应的上同调有界性，故这些满子范畴对好三角封闭。
\end{proof}

\subsection{项有界与同调有界}
\label{b4:subsec:1703:02}

“项有界”是某个具体代表的性质，“上同调有界”是导出范畴中对象的性质。上命题说明后者允许更换成项有界的代表，却没有说明该代表还能同时逐项投射或逐项内射。为获得特殊项，可能必须重新加上一段无限长但正合的尾部。

\ProseParagraphBreak
例如投射消解 \(\cdots\to P_2\to P_1\to P_0\to M\to0\)，删去增广项后，把 \(P_i\) 放在次数 \(-i\)，得到复形 \(P\)。它只有 \(H^0(P)=M\)，因此代表 \(D^b\) 的对象；但它的项可能向负次数无限延伸。把正合尾部简单截掉通常会新增上同调；正确截断需要把端点项换成核或余核，而后者未必仍投射。

\ProseParagraphBreak
这也解释了为什么两个有界条件在导出函子公式中各有用处：上同调有界用于识别对象属于哪一类导出范畴，项及其投射、内射性质则用于保证逐项计算有效。证明中不可把二者无说明地互换。

\subsection{模范畴中的基本实例}
\label{b4:subsec:1703:03}

对域 \(k\)，每个向量空间复形都能非典范地分解为零微分的上同调复形与可缩复形的直和。逐次选 \(B^n\subseteq Z^n\subseteq C^n\) 的补空间，则 \(C^n\cong B^n\oplus H^n\oplus B^{n+1}\)，微分只在最后一项到下一次数第一项之间为恒等。于是 \(K(k\text{-}\mathbf{Mod})\to D(k\text{-}\mathbf{Mod})\) 已是等价，每个拟同构都是同伦等价。这里依赖所有向量空间短正合列可分裂。

\ProseParagraphBreak
整数上情况不同。两项复形 \(P=(\mathbb Z\xrightarrow{m}\mathbb Z)\) 位于次数负一、零，其中 \(m\ge2\)，自然映到 \(\mathbb Z/m\mathbb Z\)，是拟同构而非同伦等价。它给出屋顶
\[
 \mathbb Z/m\mathbb Z\xleftarrow{\simeq}P\longrightarrow\mathbb Z[1],
\]
右边的复形映射 \(f\) 在次数负一取恒等、在次数零取零，满足链映射条件。若 \(f\) 同伦于零，同伦只能在次数零给出 \(h^0:\mathbb Z\to\mathbb Z\)，而次数负一的同伦等式要求 \(h^0\circ m=1_{\mathbb Z}\)，即某个整数乘 \(m\) 等于一，这是不可能的。\cref{b4:prop:resolution-computes-derived-hom}将证明，源为这种有上界投射复形时，局部化不再使新的同伦类变成零；因此该屋顶表示非零导出态射。\secref{b4:sec:1704}还将严格算出 \(\operatorname{Hom}_D(\mathbb Z/m\mathbb Z,\mathbb Z[1])\cong\mathbb Z/m\mathbb Z\)，说明其倍数按模 \(m\) 区分。它在所有上同调对象之间诱导零映射，因为两端非零上同调位于不同次数。这说明“知道全部上同调对象和其上诱导的映射”仍不足以恢复导出范畴的态射。

% !TeX root = ../../Book4.tex
% 纲要标题：### 18.4 消解模型
% 纲要标签：b4:sec:1704
% 纲要位置：计划纲要.md，第 2941 行。

\section{消解模型}
\label{b4:sec:1704}

有上界投射复形到零调复形的映射，以及零调复形到有下界内射复形的映射，都在同伦范畴中为零。这个消失性质使导出态射可以回到同伦类计算，并把 Ext 识别为移位 Hom。为建立单向有界模型，下面使用\cref{b4:thm:ce-resolution}的消解存在性与比较，以及\cref{b4:cor:double-complex-comparison}。


% 纲要内容（保留纲要核验项目）：
%
% * 在足够多投射对象的情形，用有界上方的投射复形处理 \(D^-\)；在足够多内射对象的情形，用有界下方的内射复形处理 \(D^+\)
% * 导出范畴中计算态射
% * Ext 与导出 Hom 的一致性
% * \(D^b\) 中的对象一般仍可能需要向一端无限延伸的投射或内射消解；不能无条件把 \(D^b\) 等同于有界投射复形的同伦范畴
%
% ---
%

\subsection{有界方向下的投射与内射复形}
\label{b4:subsec:1704:01}

\begin{lemma}\label{b4:lem:bounded-resolution-nullhomotopy}
设 \(\mathcal A\) 为局部小交换范畴，\(A\) 为任意零调复形。若 \(P\) 是逐项投射的有上界复形，则每个复形映射 \(P\to A\) 零伦；若 \(I\) 是逐项内射的有下界复形，则每个复形映射 \(A\to I\) 零伦。这里对 \(A\) 没有有界性要求。
\end{lemma}
\begin{proof}
设 \(P^n=0\) 对 \(n>N\) 成立，\(f:P\to A\)。从 \(h^{N+1}=0\) 向下构造 \(h^n:P^n\to A^{n-1}\)。已构造高次数部分时，令 \(r^n=f^n-h^{n+1}\mathrm{d}_P^n\)。由高一度的同伦等式及 \(f\) 与微分相容，得 \(\mathrm{d}_A^nr^n=0\)。因为 \(A\) 正合，\(A^{n-1}\to Z^n(A)\) 满；\(P^n\) 投射使 \(r^n\) 提升为 \(h^n\)。于是 \(f^n=\mathrm{d}_A^{n-1}h^n+h^{n+1}\mathrm{d}_P^n\)。对每个固定次数只需有限次从 \(N\) 下降，所以构造定义了全部同伦。

内射情形把范畴改为 \(\mathcal A^{\mathrm{op}}\)，把次数反向：有下界内射复形变成有上界投射复形，正合性保留。前一构造便逐度向上延拓同伦，回到 \(\mathcal A\) 后得到第二个结论。
\end{proof}

\begin{theorem}\label{b4:thm:bounded-projective-model}
设 \(\mathcal A\) 为有足够投射对象的局部小交换范畴，\(\mathcal P\) 为其投射对象满子范畴。每个有上界复形 \(X\) 都有拟同构 \(P\to X\)，其中 \(P\in K^-(\mathcal P)\)，且自然函子
\(K^-(\mathcal P)\to D^-(\mathcal A)\) 是三角等价。
\end{theorem}
\begin{theorem}\label{b4:thm:bounded-injective-model}
设 \(\mathcal A\) 为有足够内射对象的局部小交换范畴，\(\mathcal I\) 为其内射对象满子范畴。每个有下界复形 \(X\) 都有拟同构 \(X\to I\)，其中 \(I\in K^+(\mathcal I)\)，且自然函子
\(K^+(\mathcal I)\to D^+(\mathcal A)\) 是三角等价。
\end{theorem}
\begin{proof}
先证内射模型的存在。对有下界 \(X\)，应用\cref{b4:thm:ce-resolution}已经构造并比较的 Cartan--Eilenberg 内射消解，取总复形 \(I\)。若 \(X^p=0\) 对 \(p<m\) 成立，双复形可取在 \(p\ge m,q\ge0\)，所以每个总次数只含有限多个项。有限直和内射，故 \(I\) 逐项内射且有下界。增广 \(X\to I\) 是拟同构，由逐列消解的正合性与\cref{b4:cor:double-complex-comparison}得到；有限对角线使这里不涉及无限总化。投射模型则把同一构造在反范畴中对次数反向应用：所得双复形位于 \(p\le N,q\le0\)，每条对角线仍有限，给出逐项投射的 \(P\to X\)。

直接在有界方向的屋顶范畴中证明全忠实性。记 \(K^-=K^-(\mathcal A)\)，以 \(Q^-:K^-\to D^-(\mathcal A)\) 表示局部化函子，令 \(P\in K^-(\mathcal P)\)。对 \(K^-\) 中的拟同构 \(s:U\to V\)，其锥及各平移零调，故由\cref{b4:lem:bounded-resolution-nullhomotopy}及三角的 Hom 正合列，
\[
 s_*:\operatorname{Hom}_{K^-}(P,U)
 \xrightarrow{\sim}\operatorname{Hom}_{K^-}(P,V).
\]
给定 \(D^-(\mathcal A)\) 中的屋顶 \(P\xleftarrow{s}U\xrightarrow{f}Y\)，取唯一的同伦类 \(a:P\to U\)，使 \(sa=1_P\)。局部化后 \(Q^-(a)=Q^-(s)^{-1}\)，所以屋顶等于 \(Q^-(fa)\)。若 \(Q^-(g)=0\)，屋顶相等判据给出拟同构 \(t:V\to P\)，使 \(gt=0\)；取 \(b:P\to V\) 满足 \(tb=1_P\)，便得 \(g=gtb=0\)。因此对每个有上界 \(Y\)，
\[
 \operatorname{Hom}_{K^-}(P,Y)
 \xrightarrow{\sim}\operatorname{Hom}_{D^-(\mathcal A)}(P,Y).
\]

内射情形记 \(K^+=K^+(\mathcal A)\)，以 \(Q^+:K^+\to D^+(\mathcal A)\) 表示局部化函子，令 \(I\in K^+(\mathcal I)\)。同一引理的内射部分与反变 Hom 正合列说明，对每个 \(K^+\) 中的拟同构 \(s:U\to V\)，
\[
 s^*:\operatorname{Hom}_{K^+}(V,I)
 \xrightarrow{\sim}\operatorname{Hom}_{K^+}(U,I).
\]
对分母在右的屋顶 \(X\xrightarrow{f}U\xleftarrow{s}I\)，取 \(a:U\to I\) 满足 \(as=1_I\)，屋顶便由 \(af\) 表示。若 \(Q^+(g)=0\)，右屋顶的相等判据给出拟同构 \(t:I\to V\)，使 \(tg=0\)；取 \(b:V\to I\) 满足 \(bt=1_I\)，得 \(g=btg=0\)。因此对每个有下界 \(X\)，
\[
 \operatorname{Hom}_{K^+}(X,I)
 \xrightarrow{\sim}\operatorname{Hom}_{D^+(\mathcal A)}(X,I).
\]
把 \(Y\) 也取为投射模型、把 \(X\) 也取为内射模型，即得两个模型函子的全忠实性；已证的拟同构提供本质满性。由于 \(\mathcal A\) 局部小，复形映射组成集合，取同伦类后仍为集合，故上述模型也说明 \(D^-(\mathcal A)\)、\(D^+(\mathcal A)\) 在原来的集合范围内局部小。

有限直和保持投射性和内射性，所以模型范畴均对平移和映射锥封闭。上述自然函子正是局部化函子 \(Q^-\)、\(Q^+\) 在模型满子范畴上的限制，与平移相容，并把每个标准锥三角送到相应有界导出范畴的好三角。因此所得等价是三角等价。
\end{proof}

\subsection{导出范畴中态射的计算}
\label{b4:subsec:1704:02}

\begin{proposition}\label{b4:prop:resolution-computes-derived-hom}
设 \(\mathcal A\) 为局部小交换范畴，并对局部化采用\subsecref{b4:subsec:1701:03}的大小约定，简记 \(K=K(\mathcal A)\)、\(D=D(\mathcal A)\)。若 \(\mathcal A\) 中的上链复形 \(P\) 逐项投射且有上界，\(I\) 逐项内射且有下界，则对任意 \(\mathcal A\) 中的上链复形 \(X,Y\)，自然映射
\[
 \operatorname{Hom}_{K}(P,Y)\xrightarrow{\sim}\operatorname{Hom}_{D}(P,Y),
 \qquad
 \operatorname{Hom}_{K}(X,I)\xrightarrow{\sim}\operatorname{Hom}_{D}(X,I)
\]
为双射。这里不要求 \(Y\) 有上界或 \(X\) 有下界。
\end{proposition}
\begin{proof}
若 \(s:U\to V\) 是拟同构，其锥及各平移零调。前一引理给出
\[
 \operatorname{Hom}_K(P,\operatorname{Cone}(s)[n])=0\qquad(n\in\mathbb Z).
\]
三角的 Hom 正合列因而说明 \(s_*:\operatorname{Hom}_K(P,U)\to\operatorname{Hom}_K(P,V)\) 为同构。

给定屋顶 \(P\xleftarrow{s}U\xrightarrow{f}Y\)，取唯一的同伦类 \(a:P\to U\)，使 \(sa=1_P\)。局部化后 \(Q(a)=Q(s)^{-1}\)，所以该屋顶等于 \(Q(fa)\)，证明满射。若 \(Q(g)=0\)，则由屋顶判据存在拟同构 \(t:V\to P\)，使 \(gt=0\)。同样取 \(b:P\to V\) 满足 \(tb=1\)，得 \(g=gtb=0\)，证明单射。

对内射版本，使用反变 Hom 正合列、引理第二项及分母在右的屋顶。于是 \(s^*:\operatorname{Hom}_K(V,I)\to\operatorname{Hom}_K(U,I)\) 为同构；完全相同的代表与消去论证给出第二个双射。
\end{proof}

若同时有 \(P\to X\) 和 \(Y\to I\)，便可任选
\[
 \operatorname{Hom}_D(X,Y)
 \cong\operatorname{Hom}_K(P,Y)
 \cong\operatorname{Hom}_K(P,I)
 \cong\operatorname{Hom}_K(X,I)
\]
来计算。中间比较来自自然的拟同构及上命题，故不同消解不会产生互不相关的答案。

\subsection{Ext 与导出 Hom 的一致性}
\label{b4:subsec:1704:03}

\begin{theorem}\label{b4:thm:derived-ext-identification}
设 \(R\) 为含幺环，\(M,N\) 为左 \(R\)-模，视为次数零复形。则对 \(n\ge0\) 有自然同构
\[
 \operatorname{Ext}_R^n(M,N)\cong
 \operatorname{Hom}_{D(R\text{-}\mathbf{Mod})}(M,N[n]).
\]
对 \(n<0\)，右端为零；\(n=0\) 时该同构是通常的 Hom 识别。若 \(\mathcal A\) 是有足够内射对象的局部小交换范畴，\(M,N\in\mathcal A\)，则在固定的集合大小约定下也有 \(\operatorname{Ext}_{\mathcal A}^n(M,N)\cong\operatorname{Hom}_{D(\mathcal A)}(M,N[n])\)，其中 Ext 用右导出函子定义，负次数的右端同样为零。
\end{theorem}
\begin{proof}
取 \(N\to I\) 的非负次数内射消解。\(I[n]\) 仍逐项内射且有下界，所以前一命题及 Hom 复形公式给出
\[
 \operatorname{Hom}_D(M,N[n])\cong\operatorname{Hom}_K(M,I[n])
 \cong H^n\operatorname{Hom}^{\bullet}(M,I)
 =\operatorname{Ext}^n_R(M,N).
\]
因为 \(M\) 仅在零次非零，此处的 Hom 微分就是后复合 \(\mathrm{d}_I\)，与\chapref{b4:ch:08}内射定义一致。负次数的 Hom 复形项为零，所以所得群为零。比较映射唯一到同伦，故同构对两个变量自然。

也可以取 \(P\to M\) 的非正次数投射消解，得到 \(H^n\operatorname{Hom}^{\bullet}(P,N)\)。须注意\chapref{b4:ch:08}投射计算采用 \(\delta f=f\mathrm{d}_P\)，而内部 Hom 在次数 \(n\) 的微分为 \((-1)^{n+1}f\mathrm{d}_P\)。每个次数乘 \((-1)^{n(n+1)/2}\) 给出前者到后者的复形同构。连同\chapref{b4:ch:08}的平衡比较，即与上面的内射识别相容；不能只换微分而保留所有余圈代表的符号。
\end{proof}

因此 \(\operatorname{Hom}_D(\mathbb Z/m,\mathbb Z[1])\cong\mathbb Z/m\)，而 \(\operatorname{Hom}_D(\mathbb Z/m,\mathbb Z)=0\)。前一群记录扩张，后一群只记录普通同态。按照\chapref{b4:ch:09}用第一变量的内射连接同态识别扩张类的约定，一次 Yoneda 扩张对应其好三角的最后箭头 \(\delta\)。

\ProseParagraphBreak
可以直接核对这一符号。给定 \(0\to N\xrightarrow{i}E\xrightarrow{q}M\to0\)，取内射消解 \(j:N\to I\)，并延拓为 \(a:E\to I^0\)。由于 \(\mathrm{d}_Ia\) 消去 \(N\)，它唯一下降为余圈 \(f:M\to I^1\)，满足 \(fq=\mathrm{d}_Ia\)；这就是\chapref{b4:ch:09}的内射连接类。在 \(\operatorname{Cone}(i)\to I[1]\) 上，\(j[1]p_i-fr\) 的零伦取 \(h^0=a\)：次数负一的等式是 \(ai=j\)，次数零的等式是 \(-\mathrm{d}_Ia=-fq\)。所以 \(Q(j[1])\delta=Q(f)\)。若改用无符号投射余圈，必须连同前述每度重标识和平衡比较一起换算，不能据次数一的重标识孤立地改变这个扩张对应。

\subsection{有界导出对象与单向无界消解}
\label{b4:subsec:1704:04}

有界导出对象未必有有界投射代表。取域 \(k\) 及交换环 \(R=k[\varepsilon]/(\varepsilon^2)\)，令 \(k=R/(\varepsilon)\)。复形
\[
 \cdots\xrightarrow{\varepsilon}R\xrightarrow{\varepsilon}R
 \xrightarrow{\varepsilon}R\longrightarrow k\longrightarrow0
\]
是自由消解，因为 \(\ker(\varepsilon:R\to R)=\varepsilon R\)。施加 \(\operatorname{Hom}_R(-,k)\) 后所有微分为零，故
\[
 \operatorname{Ext}_R^n(k,k)\cong k\qquad(n\ge0).
\]
若 \(k\) 在导出范畴中同构于一个有界投射复形 \(Q\)，则前面的态射计算给出 \(\operatorname{Ext}_R^n(k,k)\cong\operatorname{Hom}_K(Q,k[n])\)。当 \(n\) 足够大时，\(k[n]\) 的唯一非零项已移到 \(Q\) 的范围以外，后者为零，矛盾。

\ProseParagraphBreak
所以 \(k\in D^b(R\text{-}\mathbf{Mod})\)，但不能用 \(K^b(R\text{-}\mathbf{Proj})\) 中的对象代表。与此相对，\(\mathbb Z/m\) 有长度一的自由消解，可以由有界投射复形代表。是否允许这样的有限模型，与环和模的同调维数有关；\chapref{b4:ch:20}将把有限生成投射项构成的有限模型称为完美复形。

\ProseParagraphBreak
本章的结论不包括“逐项投射的任意无界复形都能计算导出态射”。向下构造零伦时必须有起始次数；没有这个起点，逐次投射性本身不能启动归纳。无界情形所需的更强条件在\secref{b4:sec:1804}介绍。


\begin{chaptersummary}
拟同构在同伦范畴中满足交换和消去条件，因此导出态射可写成屋顶，屋顶的相等由共同细化判断。这个计算既证明局部化的泛性质，也允许把三角公理从同伦范畴传下来。由此，一般复形短正合列都能给出导出范畴中的好三角。

\ProseParagraphBreak
截断负责把上同调有界的对象换成项有界的代表，投射和内射替换负责把态射计算化为同伦计算。它们解决的问题不同：有界对象可能仍需向一端无限延伸的特殊消解。Ext 因而成为到平移对象的态射群，而这些态射又保留了普通上同调对象之间映射看不到的扩张信息。
\end{chaptersummary}

\BookExercises

\begin{exercise}\label{b4:ex:17-cyclic-ext}
设 \(m,n\ge1\)。计算 \(\operatorname{Hom}_{D(\mathbf{Ab})}(\mathbb Z/m\mathbb Z,\mathbb Z/n\mathbb Z[r])\) 对所有整数 \(r\) 的值，并写出 \(r=0,1\) 的核与余核表达。
\end{exercise}
\begin{solution}
用 \(\mathbb Z\xrightarrow{m}\mathbb Z\) 替换第一变量，Hom 复形仅在零、一次非零，两项都是 \(\mathbb Z/n\mathbb Z\)，微分为负乘 \(m\)。故零次为 \((\mathbb Z/n\mathbb Z)[m]\)，一次为 \((\mathbb Z/n\mathbb Z)/m(\mathbb Z/n\mathbb Z)\)，二者均同构于 \(\mathbb Z/\gcd(m,n)\mathbb Z\)，其余为零。零次与一次的具体子群、商群不同，不把这种同构类型误认成同一内部对象。
\end{solution}

\begin{exercise}\label{b4:ex:17-roof-localize-zero}
设 \(f:X\to Y\) 为复形映射。若它经某个零调复形分解，证明 \(Q(f)=0\)。反之，若 \(Q(f)=0\)，证明 \(f\) 在同伦范畴中经某个零调复形分解。
\end{exercise}
\begin{solution}
正向因为中间对象在 \(D\) 中为零。反向由屋顶判据取拟同构 \(s:U\to X\)，使 \(fs=0\)。补成三角 \(U\xrightarrow{s}X\to C_s\to U[1]\)。反变 Hom 正合性使 \(f\) 经 \(C_s\) 分解；拟同构的锥零调，故满足要求。
\end{solution}

\begin{exercise}\label{b4:ex:17-roof-two-denominators}
设两个屋顶 \((s,f)\)、\((s,g)\) 具有同一个左分母。证明它们在导出范畴中相等，当且仅当存在拟同构 \(u\) 使 \(fu=gu\)。
\end{exercise}
\begin{solution}
若有这样的 \(u\)，共同细化立即给出相等。反向，乘上可逆的 \(Q(s)\) 得 \(Q(f)=Q(g)\)，再用普通态射的屋顶相等判据。也可以直接从共同细化 \(sa=sb\)、\(fa=gb\) 出发，先用拟同构 \(s\) 的消去条件进一步使 \(a=b\)。
\end{solution}

\begin{exercise}\label{b4:ex:17-field-two-degrees}
设域 \(k\) 上的复形 \(X\) 仅有 \(H^{-1}(X)=k^2,H^0(X)=k\)，\(Y\) 仅有 \(H^{-1}(Y)=k,H^0(Y)=k^3\)。计算 \(\dim_k\operatorname{Hom}_D(X,Y[r])\)。所得维数是否依赖分裂的选择？
\end{exercise}
\begin{solution}
向量空间复形分解为上同调与可缩部分，故计算可在零微分复形上进行。\(r=0\) 时维数为 \(2\cdot1+1\cdot3=5\)；\(r=1\) 时从 \(X^{-1}=k^2\) 到 \(Y[1]^{-1}=k^3\)，维数为六；\(r=-1\) 时从 \(X^0=k\) 到 \(Y[-1]^0=k\)，维数为一；其余为零。分裂不典范，但 Hom 群本身及其维数由导出对象确定，不依赖选择。
\end{solution}

\begin{exercise}\label{b4:ex:17-injective-cyclic}
设整数 \(m\ge2\)。用内射消解
\[
 0\to\mathbb Z\to\mathbb Q\to\mathbb Q/\mathbb Z\to0
\]
计算 \(\operatorname{Hom}_D(\mathbb Z/m\mathbb Z,\mathbb Z[1])\)，并描述一个生成元。
\end{exercise}
\begin{solution}
内射复形取 \(I^0=\mathbb Q,I^1=\mathbb Q/\mathbb Z\)。Hom 的零次项为零，一次项为 \((\mathbb Q/\mathbb Z)[m]\)。因此所求群为 \(\mathbb Z/m\mathbb Z\)，生成元由 \(\bar1\mapsto1/m+\mathbb Z\) 给出。与投射屋顶的具体符号比较时须用正文的每度重标识，不能仅比较写下的两个数值。
\end{solution}

\begin{exercise}\label{b4:ex:17-projective-dimension}
设 \(R\) 为环，\(M\) 为左 \(R\)-模，\(d\ge0\)。证明 \(\operatorname{pd}_R M\le d\) 当且仅当对所有左模 \(N\) 都有 \(\operatorname{Hom}_D(M,N[d+1])=0\)。
\end{exercise}
\begin{solution}
由 Ext 识别，条件等价于所有 \(\operatorname{Ext}_R^{d+1}(M,N)=0\)。若有长度 \(d\) 的投射消解，消解计算直接给出消失。反向取任意投射消解，维数平移给出 \(\operatorname{Ext}_R^1(\Omega^dM,N)=0\) 对所有 \(N\) 成立，故 \(\Omega^dM\) 投射，将原消解在这里截短即可。\(d=0\) 时用一次 Ext 的投射判据。
\end{solution}

\begin{exercise}\label{b4:ex:17-cubic-truncated}
令 \(R=k[t]/(t^3)\)、\(k=R/(t)\)。构造 \(k\) 的周期投射消解，并计算它到各次移位 \(k[r]\) 的导出态射群。由此说明 \(k\) 不能由有界投射复形代表。
\end{exercise}
\begin{solution}
消解交替使用乘 \(t\) 与乘 \(t^2\)，靠近增广的箭头为乘 \(t\)。因为 \(\ker t=(t^2)\)、\(\ker t^2=(t)\)，它正合。施加 \(\operatorname{Hom}_R(-,k)\) 后微分均为零，故所求群在 \(r\ge0\) 为 \(k\)，在 \(r<0\) 为零。若存在有界投射代表，则到 \(k[r]\) 的同伦态射在充分大的 \(r\) 必为零，与计算矛盾。
\end{solution}

\begin{exercise}\label{b4:ex:17-cone-localized}
设 \(s:X\to Y\) 为拟同构，\(f:Y\to Z\) 为复形映射。证明 \(\operatorname{Cone}(fs)\) 与 \(\operatorname{Cone}(f)\) 在导出范畴中同构；这种结论在同伦范畴中是否总成立？
\end{exercise}
\begin{solution}
对严格方块 \((s,1_Z)\) 取锥映射。上同调长正合列与五引理使该锥映射拟同构，故在 \(D\) 中同构。若取零调但不可缩的 \(X\)，令 \(s:X\to0\)、\(f:0\to0\)，则两个锥分别为 \(X[1]\) 和零，在 \(K\) 中不同构。因此不能去掉局部化。
\end{solution}

\begin{exercise}\label{b4:ex:17-cohomology-detect}
设 \(R\) 为环，\(X\) 为任意左模复形。证明 \(\operatorname{Hom}_{D(R)}(R[-n],X)\cong H^n(X)\)。由此证明若所有这些群为零，则 \(X=0\)。
\end{exercise}
\begin{solution}
\(R[-n]\) 是只有一项的投射复形，可用同伦类计算导出态射。映射由 \(X^n\) 的一个余圈决定，两个映射同伦恰当其差为余边，故所求群为 \(H^n(X)\)。所有群为零表示 \(X\) 零调，等价于 \(X\) 在导出范畴中为零。
\end{solution}

\begin{exercise}\label{b4:ex:17-split-cohomology-not-split}
设整数 \(m\ge2\)，考虑由非分裂短正合列 \(0\to\mathbb Z\xrightarrow{m}\mathbb Z\to\mathbb Z/m\mathbb Z\to0\) 给出的好三角。证明其最后箭头非零，尽管它诱导的每个上同调对象之间的映射都为零。
\end{exercise}
\begin{solution}
若最后箭头为零，由分裂三角判据，\(\mathbb Z\to\mathbb Z/m\mathbb Z\) 在导出范畴中有截面。取 \(H^0\) 得交换群满射的截面，与 \(\mathbb Z\) 无有限阶元素矛盾。所以最后箭头非零。它从 \(\mathbb Z/m\mathbb Z\) 指向 \(\mathbb Z[1]\)，两端非零上同调的次数不同，因此每个 \(H^n\) 上的映射均零。
\end{solution}

\begin{exercise}\label{b4:ex:17-quasi-between-models}
设 \(P,Q\) 为逐项投射且有上界的复形，\(s:P\to Q\) 为拟同构。证明它已经是同伦等价，并说明不能仅以“拟同构的锥正合”结束证明。
\end{exercise}
\begin{solution}
锥 \(C_s\) 也逐项投射且有上界。把零伦引理应用于其恒等映射 \(C_s\to C_s\)，因目标正合，得到锥可缩。三角 Hom 正合列说明 \(s\) 在 \(K\) 中同构，即同伦等价。正合本身不足以推出可缩，这里必须同时使用锥逐项投射及有上界。
\end{solution}

\begin{exercise}\label{b4:ex:17-truncation-object}
设复形 \(X\) 只有次数 \(n\) 的上同调可能非零。构造 \(D(\mathcal A)\) 中 \(X\cong H^n(X)[-n]\) 的同构，不选各项短正合列的分裂。
\end{exercise}
\begin{solution}
取 \(\tau_{\le n}X\to X\)，它是拟同构。再将 \(\tau_{\le n}X\) 的次数 \(n\) 映到 \(H^n(X)\)，其他次数映到零，得到复形映射 \(\tau_{\le n}X\to H^n(X)[-n]\)。它在唯一可能非零的上同调上为恒等，因而也是拟同构。这两个箭头给出所需同构；全部映射只用核和商，无须分裂。
\end{solution}

\begin{exercise}\label{b4:ex:17-localization-functor}
设 \(F:\mathcal A\to\mathcal B\) 为交换范畴间的正合加性函子。证明逐项 \(F\) 诱导三角函子 \(D(\mathcal A)\to D(\mathcal B)\)，并说明为何仅加性不足。
\end{exercise}
\begin{solution}
正合性使 \(H^n(FC)\cong F(H^nC)\)，故逐项 \(F\) 保持拟同构，由局部化泛性质下降。加性使它保持有限双积和平移的符号，也保持锥矩阵，所以保持好三角。仅加性时拟同构可能不再保持，例如把 \(\mathbb Z\xrightarrow{2}\mathbb Z\to\mathbb Z/2\mathbb Z\) 的投射替换与 \(\mathbb Z/2\mathbb Z\) 张量会新增负一次上同调，因而不能直接下降。
\end{solution}

\begin{exercise}\label{b4:ex:17-abelian-embedding}
设 \(\mathcal A\) 为有足够内射对象的局部小交换范畴。证明将对象放在次数零给出全忠实函子 \(\mathcal A\to D^b(\mathcal A)\)，但这不使 \(D^b(\mathcal A)\) 自动成为交换范畴。
\end{exercise}
\begin{solution}
用有下界内射消解在 \(D^+(\mathcal A)\) 中计算，次数零给出 \(\operatorname{Hom}_{D^+}(M,N)=\operatorname{Hom}_{\mathcal A}(M,N)\)。有界嵌入的全忠实性使它也等于 \(D^b\) 中的 Hom，因此对象放在次数零的函子全忠实。

若 \(D^b(\mathbf{Ab})\) 同时为交换范畴，则任意态射 \(f:X\to Y\) 可分解为满态射 \(e:X\to\operatorname{im}f\) 和单态射 \(m:\operatorname{im}f\to Y\)。由\cref{b4:prop:triangle-mono-splits}，可取 \(a:\operatorname{im}f\to X\)、\(b:Y\to\operatorname{im}f\)，满足 \(ea=1\)、\(bm=1\)。于是 \(t=ab:Y\to X\) 满足
\[
 ftf=(me)(ab)(me)=m(ea)(bm)e=me=f.
\]
对次数零的乘二态射 \(\mathbb Z\to\mathbb Z\)，全忠实性使 \(t\) 仍由一个整数表示，因而要求 \(4t=2\)，矛盾。可见原交换范畴作为满子范畴嵌入，并不使核余核的性质在较大范畴中自动成立。
\end{solution}
